\documentclass[12pt,a4paper]{article}

\usepackage[UTF8,fontset=fandol]{ctex}
\usepackage[margin=1in]{geometry}
\usepackage{amsmath,amssymb,amsthm}
\usepackage{graphicx}
\usepackage{booktabs}
\usepackage{array}
\usepackage{enumitem}
\usepackage{hyperref}
\usepackage{xcolor}
\usepackage{caption}
\usepackage{tcolorbox}
\usepackage{multirow}
\usepackage{makecell}

% === Proposal Colors ===
\definecolor{propOne}{HTML}{C0392B}
\definecolor{propTwo}{HTML}{2980B9}
\definecolor{propThree}{HTML}{16A085}
\definecolor{propFour}{HTML}{8E44AD}
\definecolor{propFive}{HTML}{D35400}
\definecolor{accent}{HTML}{2C3E50}

% === Proposal Box ===
\newtcolorbox{proposalbox}[2][]{%
  colback=#2!5,colframe=#2!70,boxrule=0.5pt,
  fonttitle=\bfseries,title={#1},
  left=6pt,right=6pt,top=4pt,bottom=4pt
}

% === Insight Box ===
\newtcolorbox{insightbox}[1][]{%
  colback=yellow!5,colframe=yellow!60!black,boxrule=0.4pt,
  fonttitle=\bfseries,title={Core Insight},#1
}

% === Risk Box ===
\newtcolorbox{riskbox}[1][]{%
  colback=red!3,colframe=red!40,boxrule=0.3pt,
  fonttitle=\bfseries,title={Open Risks},#1
}

% === Experiment Box ===
\newtcolorbox{expbox}[1][]{%
  colback=blue!3,colframe=blue!40,boxrule=0.3pt,
  fonttitle=\bfseries,title={Key Experiments},#1
}

% === Why Box ===
\newtcolorbox{whybox}[1][]{%
  colback=gray!5,colframe=gray!60,boxrule=0.4pt,
  fonttitle=\bfseries,title={First-Principles Chain},#1
}

\hypersetup{
  colorlinks=true,
  linkcolor=black!70,
  citecolor=blue!60!black,
  urlcolor=blue!70!black,
}

\title{\textbf{VLM 多模态 High-Impact Project Proposals}\\
\large{从第一性原理出发的五个研究方向}\\[6pt]
\normalsize{Independent Research Directions for Vision-Language Models}}

\author{Jinhong Lin}
\date{March 2026}

\begin{document}

\maketitle

\section*{概览}

本文档提出 5 个彼此区分明确的 VLM 研究方向。
筛选标准只有三个：
\begin{enumerate}[leftmargin=2em]
  \item 不是只靠 scaling 就能自然解决的问题。
  \item 如果做成，会改变 VLM 的能力边界，而不只是刷高单一 benchmark。
  \item 能从第一性原理往下追问，定位到结构性 root cause。
\end{enumerate}

\begin{insightbox}
如果连续追问 ``why''，当前 VLM 的主要瓶颈并不是``模型还不够大''，
而是 5 类结构性错配：
\begin{enumerate}[leftmargin=2em]
  \item 它把视频当作一堆帧，而不是随时间演化的状态。
  \item 它把空间看成 2D pattern matching，而不是可操作的 3D geometry。
  \item 它把外部知识当作附加文本，而不是与视觉证据共同裁决答案的证据系统。
  \item 它把长上下文 demonstration 当作 prompt stuffing，而不是临时可编译的 task adaptation。
  \item 它的 benchmark 仍然允许 shortcut，所以整个领域可能在优化``看起来会看图''，而不是``真的依赖图像推理''。
\end{enumerate}
\end{insightbox}

\subsection*{快速优先级}

\renewcommand{\arraystretch}{1.25}
\begin{center}
\begin{tabular}{>{\raggedright\arraybackslash}p{3.3cm} >{\raggedright\arraybackslash}p{4.8cm} >{\centering\arraybackslash}p{1.2cm} >{\centering\arraybackslash}p{1.2cm} >{\raggedright\arraybackslash}p{3.3cm}}
\toprule
\textbf{方向} & \textbf{Root Cause} & \textbf{影响力} & \textbf{风险} & \textbf{适合的人} \\
\midrule
Stateful Video VLM & 视频被压缩成一次性 token & 很高 & 高 & long-video / agent / world understanding \\
Geometry-Grounded VLM & 缺少 metric 级空间表征 & 很高 & 高 & robotics / 3D / embodied AI \\
Retrieval-Calibrated Expert VLM & 外部知识与视觉证据缺少仲裁机制 & 很高 & 中 & medical / science / deployment \\
Context-to-Adapter VLM & many-shot context 不能转化为临时适配 & 高 & 中 & efficient adaptation / open ecosystem \\
Intervention-Centric Benchmark & 评测奖励 shortcut 而非真理解 & 很高 & 中 & benchmark + method + field shaping \\
\bottomrule
\end{tabular}
\end{center}

\newpage
\section{Proposal 1: Stateful Video VLM}

\begin{proposalbox}[从``看很多帧''转向``维护可推理的事件状态'']{propOne}
\textbf{一句话摘要：}
长视频理解失败，不是因为模型没看够帧，
而是因为它没有形成一个可以持续更新、回溯、查询的事件状态系统。
\end{proposalbox}

\subsection{第一性原理推导}

\begin{whybox}
\textbf{Why} are long videos hard?\\
$\rightarrow$ 因为 token budget 固定，而视频信息量随时长近似线性增长。\\[3pt]
\textbf{Why} doesn't a longer context solve this?\\
$\rightarrow$ 因为注意力会被大量低价值帧稀释，关键因果事件反而更难保留。\\[3pt]
\textbf{Why} does naive compression still fail?\\
$\rightarrow$ 因为当前压缩大多保留``看到了什么''，却丢掉``什么时候发生、因谁而起、后面会影响什么''。\\[3pt]
\textbf{Root cause:}\\
现有 video VLM 本质上仍把视频当成 frame bag，而不是 state transition process。
\end{whybox}

\subsection{核心假设}

如果让 VLM 学会显式维护 \texttt{event state memory}，
并把后续推理建立在``事件状态的演化''而不是``帧级摘要''上，
那么它在长视频上的性能会比单纯扩 context 更稳、更便宜，也更可解释。

\subsection{方法草案}

\begin{enumerate}[leftmargin=2em]
  \item 设计 \texttt{state slots}，每个 slot 表示一个正在演化的事件或实体状态。
  \item 引入三类动作：\texttt{write}、\texttt{update}、\texttt{rewind}。
  \item 构造 ``same frames, different order'' 的反事实视频，让模型必须学习顺序与因果。
  \item 对 memory 做 question-conditioned credit assignment，逼迫模型只保留未来问答真正需要的状态。
  \item 先用小模型验证 memory 结构本身带来的增益，再迁移到大模型。
\end{enumerate}

\begin{expbox}
\begin{enumerate}[leftmargin=2em]
  \item 主 benchmark 用 \href{https://arxiv.org/abs/2405.21075}{Video-MME} 和 \href{https://arxiv.org/abs/2406.04264}{MLVU}。
  \item 重点看长视频分桶，不只看 overall score。
  \item 增加 order-sensitive 和 causal-sensitive stress test。
  \item 额外评估 memory trace 是否可解释：错误答案时，memory 中是否早已丢掉关键事件。
\end{enumerate}
\end{expbox}

\subsection{为什么 high impact}

长视频是 agent、监控、手术录像、教育录像、会议分析、机器人长期任务理解的共同底层能力。
如果这个方向做出来，它不会只影响一个 benchmark，而会影响所有需要在时间上持续建模世界的 VLM 系统。

\begin{riskbox}
风险在于 memory 机制可能只是在架构上增加复杂度，却没有真的学到 state abstraction。\\
\textbf{Mitigation:}
先做两个最小闭环：
\begin{enumerate}[leftmargin=2em]
  \item 用合成视频验证模型是否真的学会``顺序敏感''，而不是更强 OCR。
  \item 用可控小游戏环境生成视频，让 ground-truth state 可直接监督。
\end{enumerate}
\end{riskbox}

\newpage
\section{Proposal 2: Geometry-Grounded VLM}

\begin{proposalbox}[让 VLM 从``看起来懂空间''变成``真的会算空间'']{propTwo}
\textbf{一句话摘要：}
很多 VLM 在图像问答里像是会空间推理，
但它们真正缺的是 metric-aware 的几何表征，
所以一旦任务需要可执行的 3D 推理，性能会快速坍塌。
\end{proposalbox}

\subsection{第一性原理推导}

\begin{whybox}
\textbf{Why} do current VLMs look strong on many visual tasks?\\
$\rightarrow$ 因为大量 benchmark 只要求语义识别，不要求可执行的空间量化。\\[3pt]
\textbf{Why} do they fail in robotics, AR, and navigation?\\
$\rightarrow$ 因为这些任务需要距离、深度、遮挡、可达性、相对尺寸等可计算关系。\\[3pt]
\textbf{Why} can't language CoT fix this directly?\\
$\rightarrow$ 因为空间错误主要不是 reasoning format 问题，而是输入表征里根本没有稳定几何量。\\[3pt]
\textbf{Root cause:}\\
当前 VLM 学到的是 image-text correlation，不是 geometry-faithful world representation。
\end{whybox}

\subsection{核心假设}

如果把 3D geometry 明确提升为中间表征，并让语言推理建立在几何对象与关系图之上，
VLM 才可能真正获得可转移的 spatial reasoning 能力。

\subsection{方法草案}

\begin{enumerate}[leftmargin=2em]
  \item 从图像或多视角输入中构建轻量级 scene graph、point tokens、object-centric geometry tokens。
  \item 所有空间问题都先转换成对几何对象的引用，再进行语言推理。
  \item 构造几何反事实数据，例如保持语义基本不变，只改距离、朝向、遮挡或大小关系。
  \item 对输出施加 \texttt{geometry consistency loss}，要求答案与中间几何表征一致。
\end{enumerate}

\begin{expbox}
\begin{enumerate}[leftmargin=2em]
  \item 基础评估可参考 \href{https://arxiv.org/abs/2401.12168}{SpatialVLM} 一类的 3D spatial reasoning 设定。
  \item 在合成场景中测量距离估计误差、相对大小排序错误率、遮挡判断准确率。
  \item 在下游场景验证，例如机器人 pick-and-place、室内导航、AR 指令理解。
\end{enumerate}
\end{expbox}

\subsection{为什么 high impact}

一旦空间推理从``描述层''走到``执行层''，
VLM 才真正连接到 embodied AI、机器人操作、工业检测和真实世界代理。
这是从 demo intelligence 走向 operational intelligence 的关键一步。

\begin{riskbox}
几何中间层可能带来额外噪声，尤其是单图 3D 恢复不稳定时。\\
\textbf{Mitigation:}
先做 object-centric 相对几何，再逐步过渡到绝对 metric geometry；先证明关系正确，再追求尺度也准。
\end{riskbox}

\newpage
\section{Proposal 3: Retrieval-Calibrated Expert VLM}

\begin{proposalbox}[面向医学/科学的证据仲裁型多模态系统]{propThree}
\textbf{一句话摘要：}
专家场景下，真正难的不是``图像 + 文本一起输入''，
而是视觉证据、外部知识、任务规范三者经常冲突，
现有 VLM 缺少一个显式的证据仲裁机制。
\end{proposalbox}

\subsection{第一性原理推导}

\begin{whybox}
\textbf{Why} do general-purpose VLMs drop sharply in expert domains?\\
$\rightarrow$ 因为这些任务的答案通常不能只从图像里直接读出。\\[3pt]
\textbf{Why} is domain fine-tuning alone insufficient?\\
$\rightarrow$ 因为专家知识会变化、会版本化，而且很多关键规则本来就应该留在参数外部。\\[3pt]
\textbf{Why} is multimodal RAG harder than text RAG?\\
$\rightarrow$ 因为视觉所见与检索到的文本证据可能冲突，而现有模型通常没有明确地区分``我看到了什么''和``外部资料说了什么''。\\[3pt]
\textbf{Root cause:}\\
当前 VLM 没有把答案看作一个由多源证据共同裁决的过程，而是在做 end-to-end next-token imitation。
\end{whybox}

\subsection{核心假设}

如果把视觉发现、检索证据、任务规范拆成三个显式证据层，
并训练模型进行跨证据一致性判断与 selective answering，
那么专家型多模态系统会显著更可靠。

\subsection{方法草案}

\begin{enumerate}[leftmargin=2em]
  \item 模型先生成 \texttt{visual findings}，明确区分可见事实和不可见推断。
  \item Retriever 只根据 findings 和 query 去抓取规范、论文、病例知识库。
  \item 构建 \texttt{evidence graph}，边表示支持、冲突、缺失、需要补充证据。
  \item 引入 \texttt{answer / abstain / request-more-context} 三路决策。
  \item 在训练中加入冲突证据对，让模型学会在视觉证据与文本先验矛盾时优先报告不确定性。
\end{enumerate}

\begin{expbox}
\begin{enumerate}[leftmargin=2em]
  \item 用 \href{https://arxiv.org/abs/2406.19593}{SK-VQA} 这类 context-augmented multimodal setting 做基础评估。
  \item 医疗场景可参考 \href{https://arxiv.org/abs/2405.03162}{Med-Gemini} 暴露出来的问题，尤其是 3D CT 报告和临床可接受性。
  \item 除了最终准确率，还要评估 evidence attribution、citation faithfulness、校准误差和 abstention quality。
\end{enumerate}
\end{expbox}

\subsection{为什么 high impact}

这是最接近真实落地价值的 VLM 方向之一。
医疗、科学助手、专利检索、工业检修、法律图文材料审阅，
本质上都不是``识图''，而是``视觉证据参与知识密集型判断''。

\begin{riskbox}
风险在于系统最后退化为更复杂的 RAG pipeline，但不产生新的学习机制。\\
\textbf{Mitigation:}
关键不是只搭系统，而是把``证据冲突仲裁''做成新的训练目标和 benchmark，
逼迫模型学习何时相信图像、何时相信检索、何时必须 abstain。
\end{riskbox}

\newpage
\section{Proposal 4: Context-to-Adapter VLM}

\begin{proposalbox}[把 many-shot multimodal ICL 变成临时适配能力]{propFour}
\textbf{一句话摘要：}
many-shot multimodal ICL 说明长上下文里藏着巨大的适配潜力，
但现在的用法太粗糙了。
真正的机会不是塞更多示例，而是把示例``编译''为临时任务适配器。
\end{proposalbox}

\subsection{第一性原理推导}

\begin{whybox}
\textbf{Why} is many-shot ICL important?\\
$\rightarrow$ 因为它证明模型在不更新参数的情况下，仍然能从大量示例里获得显著增益。\\[3pt]
\textbf{Why} are current methods still inefficient?\\
$\rightarrow$ 因为 demonstration 只是被串接进 prompt，没有被压缩成持久、可复用、任务相关的内部状态。\\[3pt]
\textbf{Why} does this become a bottleneck?\\
$\rightarrow$ 因为长 prompt 成本高、延迟高，而且示例一换，模型就得重新读一遍教材。\\[3pt]
\textbf{Root cause:}\\
当前 ICL 仍是 retrieval-and-read，而不是 retrieval-and-compile。
\end{whybox}

\subsection{核心假设}

如果把 many-shot multimodal context 显式蒸馏成一个临时 adapter、latent task sketch 或 routing state，
模型就能更便宜、更稳定地适应新 domain，同时保留 ICL 的免微调优势。

\subsection{方法草案}

\begin{enumerate}[leftmargin=2em]
  \item 先检索大量 demonstration，但不直接把它们完整送入主模型。
  \item 用一个 \texttt{context compiler} 把这些示例压缩成临时 task adapter，例如 prefix state、LoRA-like ephemeral weights 或 latent task program。
  \item 主模型只消费``编译结果 + 少量代表性样例''，而不是完整原始示例。
  \item 训练目标直接优化给定大量示例后的 adaptation quality，而不是普通 next-token loss。
  \item 加入 uncertainty-aware example selection，减少无关示例对临时适配器的污染。
\end{enumerate}

\begin{expbox}
\begin{enumerate}[leftmargin=2em]
  \item 以 \href{https://arxiv.org/abs/2405.09798}{Many-Shot In-Context Learning in Multimodal Foundation Models} 为出发点，关注从 few-shot 到近 2000 例的 scaling 区间。
  \item 比较三种成本曲线：原始 many-shot ICL、full fine-tuning、context-to-adapter。
  \item 任务覆盖分类、VQA、定位、医疗影像、遥感、分子图像等跨域设置。
  \item 重点汇报 accuracy-latency-cost 三维 trade-off，而不是只报一个精度数。
\end{enumerate}
\end{expbox}

\subsection{为什么 high impact}

这条线最有机会改变多模态模型的部署范式。
它指向一个现实命题：能否让 VLM 在不重新训练大模型的前提下，
快速适配一个新任务、新医院、新工厂、新卫星数据源。

\begin{riskbox}
临时 adapter 可能学到的只是 prompt compression，而不是 task adaptation。\\
\textbf{Mitigation:}
做强分布外测试：训练时给的是某类任务 demonstration，
测试时要求迁移到结构相近但表面形式不同的新任务，
验证其是否真的学到了可迁移 task program。
\end{riskbox}

\newpage
\section{Proposal 5: Intervention-Centric Benchmark and Training}

\begin{proposalbox}[重新定义``真的会看图'']{propFive}
\textbf{一句话摘要：}
如果 benchmark 本身奖励 shortcut，那么整个领域都会被带偏。
高影响的一步不是再堆一个更大的模型，
而是发明一套让``真假多模态能力''彻底分家的评测与训练框架。
\end{proposalbox}

\subsection{第一性原理推导}

\begin{whybox}
\textbf{Why} do many VLM scores look strong but break in real use?\\
$\rightarrow$ 因为 benchmark 往往允许语言先验、OCR shortcut、选项猜测、数据泄漏等旁门左道。\\[3pt]
\textbf{Why} is this worse than one model failing?\\
$\rightarrow$ 因为 benchmark 决定整个社区在优化什么。\\[3pt]
\textbf{Why} isn't cleaning the dataset enough?\\
$\rightarrow$ 因为只要没有 intervention，模型仍可能通过统计相关性而不是因果依赖来做题。\\[3pt]
\textbf{Root cause:}\\
当前领域对``模型到底因为什么答对''缺乏强识别力。
\end{whybox}

\subsection{核心假设}

只有当 benchmark 包含系统性的图像干预、文本干预、检索干预、选项干预，
并能明确测出``答案是否随真正因果因素而改变''时，
我们才能开始真正训练 grounded VLM。

\subsection{方法草案}

\begin{enumerate}[leftmargin=2em]
  \item 以 \href{https://arxiv.org/abs/2311.16502}{MMMU} 为 base setting，但进一步采用 \href{https://arxiv.org/abs/2409.02813}{MMMU-Pro} 那种更严格的思路，系统消除 text-only shortcut。
  \item 为每个样本构造至少一组 intervention pair：图像改了但文本不改；文本改了但图像不改；检索上下文改了但核心视觉事实不改；选项扰动使猜测失效。
  \item 指标不只看 accuracy，还看 \texttt{intervention consistency}、\texttt{causal sensitivity}、\texttt{spurious invariance}。
  \item 在 benchmark 之上再做训练方法：对同一问题的 intervention pair 施加对比式 grounded training，让模型只对真正因果证据敏感。
\end{enumerate}

\begin{expbox}
\begin{enumerate}[leftmargin=2em]
  \item 先在 MMMU / MMMU-Pro 上复现实证，确认现有模型的 shortcut 类型。
  \item 再构建新的 intervention split，并测不同模型在因果敏感性上的排名是否发生重排。
  \item 用 \href{https://arxiv.org/abs/2406.11839}{mDPO} 暴露出的 unconditional preference problem 作为训练侧动机，
  测试新训练目标能否减少``忽略图像条件''的现象。
\end{enumerate}
\end{expbox}

\subsection{为什么 high impact}

这类工作有明显的 field-shaping potential。
它不仅能发 benchmark paper，还能影响后续 alignment、data curation、RL、RAG、reasoning 的评价标准。
高影响论文很多时候不是``最强模型''，而是``改写了别人怎么证明自己强''。

\begin{riskbox}
如果只停留在 benchmark，容易被看成评测工程。\\
\textbf{Mitigation:}
必须把 benchmark 和训练方法绑定起来，形成一个闭环：
先用 intervention 暴露 shortcut，再用 intervention-aware training 消掉 shortcut。
\end{riskbox}

\section*{推荐切入点}

如果目标是``high impact + 能讲出强 why-chain''，我最推荐三个切入点：
\begin{enumerate}[leftmargin=2em]
  \item \textbf{Proposal 1}：适合 long-video / agent / world understanding，大问题、大空间，但工程难度高。
  \item \textbf{Proposal 3}：最接近现实价值，最容易讲 application story，尤其适合医疗、科学、工业知识系统。
  \item \textbf{Proposal 5}：最有可能做出 field-shaping 工作，尤其适合 benchmark + method + theory positioning 路线。
\end{enumerate}

\section*{参考依据}

\begin{enumerate}[leftmargin=2em]
  \item \href{https://arxiv.org/abs/2311.16502}{MMMU}.
  \item \href{https://arxiv.org/abs/2409.02813}{MMMU-Pro}.
  \item \href{https://arxiv.org/abs/2405.21075}{Video-MME}.
  \item \href{https://arxiv.org/abs/2406.04264}{MLVU}.
  \item \href{https://arxiv.org/abs/2401.12168}{SpatialVLM}.
  \item \href{https://arxiv.org/abs/2405.09798}{Many-Shot In-Context Learning in Multimodal Foundation Models}.
  \item \href{https://arxiv.org/abs/2406.19593}{SK-VQA}.
  \item \href{https://arxiv.org/abs/2405.03162}{Advancing Multimodal Medical Capabilities of Gemini}.
  \item \href{https://arxiv.org/abs/2406.11839}{mDPO}.
\end{enumerate}

\end{document}
